\documentclass[11pt]{article}
\usepackage[utf8]{inputenc}
\usepackage[margin=1in]{geometry}
\usepackage{amsmath}
\usepackage{graphicx}
\usepackage{booktabs}
\usepackage{hyperref}
\usepackage[round]{natbib}
\usepackage{lineno}
\usepackage{xcolor}
\usepackage{multirow}
\usepackage{float}

\linenumbers

\title{EEG-fMRI in Awake Rats Reveals Suppressed Whole-Brain Sensory Responsiveness During Absence Seizures, Corroborated by Mean-Field Simulations}

\author{
Author A$^{1,*}$, Author B$^{1}$, Author C$^{2}$, Author D$^{1}$, Author E$^{1}$ \\
\\
$^{1}$Grenoble Institut des Neurosciences, Universit\'{e} Grenoble Alpes, Inserm \\
$^{2}$Centre de Recherche en Neurosciences de Lyon, CNRS \\
$^{*}$Corresponding author: authora@univ-grenoble-alpes.fr
}

\date{}

\begin{document}
\maketitle

\begin{abstract}
Absence seizures cause transient impairment of consciousness, yet how the brain processes incoming sensory information during spike-and-wave discharges (SWDs) remains unknown because prior human neuroimaging studies were conducted exclusively at rest. Here we established a zero-echo-time (ZTE) fMRI protocol for awake, non-curarized GAERS rats (11 animals, 6 male/5 female), achieving 35.8~dB lower acoustic noise than EPI (78.7 vs. 114.5~dB) and mean framewise displacement of only $0.69 \pm 0.32$~$\mu$m. During simultaneous EEG-fMRI, visual and whisker stimulation were delivered during both interictal and ictal (spontaneous SWD) periods. Whole-brain hemodynamic responses to both stimulus modalities were globally suppressed and spatially restricted during ictal compared to interictal periods (F-contrast, $p < 0.05$, cluster-corrected). Cortical hemodynamic responses became negatively polarized during seizures despite ongoing stimulation, with mean BOLD signal change of $-0.18 \pm 0.06$\% (ictal) vs. $+0.31 \pm 0.08$\% (interictal) in visual cortex ($t(10) = 4.82$, $p < 0.001$, paired $t$-test). A mean-field simulation using an Adaptive Exponential (AdEx) neuron model reproduced SWD dynamics and demonstrated restricted propagation of stimulation-evoked activity during the ictal regime, matching fMRI observations. These findings provide the first direct evidence that absence seizures actively suppress sensory processing across the entire brain.
\end{abstract}

\section{Introduction}

Absence epilepsy is a common pediatric epilepsy syndrome characterized by brief episodes of impaired consciousness associated with bilateral synchronous spike-and-wave discharges (SWDs) at 2--5~Hz \citep{panayiotopoulos2008typical}. Beyond the transient consciousness impairments during seizures, patients often exhibit neuropsychiatric comorbidities including deficits in attention, memory, and mood \citep{blumenfeld2005consciousness, mccarthy2015attention}. Despite extensive clinical characterization, the neural mechanism by which SWDs impair responsiveness to environmental stimuli remains poorly understood.

Human neuroimaging studies of absence epilepsy have been limited to resting-state paradigms, revealing widespread cortical deactivation and thalamic activation during seizures \citep{tenney2014seizures}. However, no study has examined how the brain processes externally delivered sensory stimulation during seizures. Rodent fMRI studies in genetic models such as GAERS \citep{vergnes1982spontaneous, depaulis2015gaers} have provided valuable insights into seizure hemodynamics, but these typically employed curarized or anesthetized preparations, confounding interpretation.

A major technical barrier to awake-state fMRI in rodents is the loud acoustic noise generated by standard echo-planar imaging (EPI) sequences, which causes stress, increases motion artifacts, and suppresses spontaneous seizure generation in epilepsy models. Zero-echo-time (ZTE) sequences offer a solution by dramatically reducing acoustic noise while maintaining sensitivity to hemodynamic changes through T1-relaxation rate mechanisms rather than T2*-weighted contrast \citep{madsen2015zte}.

Here we combine awake-state ZTE fMRI with simultaneous EEG recording and controlled sensory stimulation in GAERS rats, enabling the first direct comparison of brain-wide sensory responses during ictal versus interictal states. We complement the imaging data with mean-field computational simulations that reproduce the observed suppression pattern.

\section{Results}

\subsection{ZTE fMRI Enables Awake Imaging with Minimal Motion}

The ZTE sequence achieved substantially lower acoustic noise and motion compared to standard EPI (Table~\ref{tab:zte_params}).

\begin{table}[H]
\centering
\caption{Comparison of ZTE and EPI acquisition parameters and motion metrics. Values are mean $\pm$ SD across all included sessions ($n = 11$ animals). Framewise displacement computed following \citet{power2012spurious}. Motion occurrences defined as frames exceeding 50~$\mu$m displacement.}
\label{tab:zte_params}
\begin{tabular}{lccc}
\toprule
\textbf{Parameter} & \textbf{ZTE (this study)} & \textbf{EPI (prior studies)} & \textbf{$p$-value} \\
\midrule
Peak acoustic noise (dB) & $78.7$ & $114.5$ & --- \\
Noise reduction & $35.8$ dB ($\sim$62$\times$ lower) & Reference & --- \\
Mean framewise displacement ($\mu$m) & $0.69 \pm 0.32$ & --- & --- \\
Mean head rotation (deg) & $0.79 \pm 0.55$ & --- & --- \\
Max displacement ($\mu$m) & $12.8 \pm 11.6$ & --- & --- \\
Max displacement (voxels) & $0.03 \pm 0.02$ & --- & --- \\
Motion occurrences/min & $0.43 \pm 0.45$ & $1.00 \pm 0.20$ & $< 0.01$ \\
Signal change (visual stim, \%) & $0.31 \pm 0.08$ & $1.2 \pm 0.3$ & --- \\
Spatial SNR & $\sim$25 & --- & --- \\
Temporal SNR & $42 \pm 12$ & --- & --- \\
\bottomrule
\end{tabular}
\end{table}

\subsection{Sensory Responses Are Globally Suppressed During Seizures}

During interictal periods, visual stimulation evoked robust positive BOLD responses throughout the visual pathway. During ictal periods, these responses were suppressed and often reversed in polarity (Table~\ref{tab:visual}).

\begin{table}[H]
\centering
\caption{Visual stimulation responses in key brain regions during interictal vs. ictal states. BOLD signal change (\%) extracted from ROIs. Statistical comparison by paired $t$-test across animals ($n = 11$). Significance assessed at $p < 0.05$ with Bonferroni correction for 6 ROIs.}
\label{tab:visual}
\begin{tabular}{lcccccc}
\toprule
\textbf{Brain Region} & \textbf{Interictal (\%)} & \textbf{Ictal (\%)} & \textbf{$\Delta$ (\%)} & \textbf{$t$-stat} & \textbf{df} & \textbf{$p$-value} \\
\midrule
Visual cortex (bilat.) & $+0.31 \pm 0.08$ & $-0.18 \pm 0.06$ & $-0.49$ & $4.82$ & 10 & $< 0.001$ \\
Superior colliculus & $+0.24 \pm 0.07$ & $-0.08 \pm 0.05$ & $-0.32$ & $3.64$ & 10 & $< 0.01$ \\
Lateral geniculate (LGN) & $+0.19 \pm 0.06$ & $-0.05 \pm 0.04$ & $-0.24$ & $3.21$ & 10 & $< 0.01$ \\
Prelimbic cortex & $+0.16 \pm 0.05$ & $-0.11 \pm 0.05$ & $-0.27$ & $3.38$ & 10 & $< 0.01$ \\
Cingulate cortex & $+0.14 \pm 0.05$ & $-0.09 \pm 0.04$ & $-0.23$ & $3.12$ & 10 & $< 0.05$ \\
Secondary motor cortex & $+0.12 \pm 0.04$ & $-0.06 \pm 0.03$ & $-0.18$ & $2.88$ & 10 & $< 0.05$ \\
\bottomrule
\end{tabular}
\end{table}

Whisker stimulation showed a parallel pattern of suppression across the somatosensory pathway (Table~\ref{tab:whisker}).

\begin{table}[H]
\centering
\caption{Whisker stimulation responses during interictal vs. ictal states. Same conventions as Table~\ref{tab:visual}. Bonferroni correction for 3 ROIs.}
\label{tab:whisker}
\begin{tabular}{lcccccc}
\toprule
\textbf{Brain Region} & \textbf{Interictal (\%)} & \textbf{Ictal (\%)} & \textbf{$\Delta$ (\%)} & \textbf{$t$-stat} & \textbf{df} & \textbf{$p$-value} \\
\midrule
Barrel field S1 & $+0.28 \pm 0.07$ & $-0.14 \pm 0.05$ & $-0.42$ & $4.38$ & 10 & $< 0.001$ \\
VPM thalamus & $+0.21 \pm 0.06$ & $-0.06 \pm 0.04$ & $-0.27$ & $3.52$ & 10 & $< 0.01$ \\
Frontal cortex & $+0.15 \pm 0.05$ & $-0.10 \pm 0.04$ & $-0.25$ & $3.28$ & 10 & $< 0.01$ \\
\bottomrule
\end{tabular}
\end{table}

F-contrast maps confirmed significant interictal-minus-ictal differences across both stimulus modalities ($p < 0.05$, cluster-corrected). The suppression was not limited to primary sensory cortex but extended to thalamic relays, frontal attentional regions, and subcortical structures.

\subsection{Seizure Hemodynamic Signature}

Analysis of seizure-only epochs (without stimulation) revealed characteristic hemodynamic patterns across cortical and subcortical regions (Table~\ref{tab:seizure_hemo}).

\begin{table}[H]
\centering
\caption{Hemodynamic responses during seizures (without stimulation). F-contrast from GLM with seizure regressor. Polarity indicates direction of BOLD signal change during SWDs. Values are mean $\pm$ SD across $n = 11$ animals. Cluster-corrected at $p < 0.05$.}
\label{tab:seizure_hemo}
\begin{tabular}{lcccc}
\toprule
\textbf{Brain Region} & \textbf{Polarity} & \textbf{BOLD Change (\%)} & \textbf{$F$-stat} & \textbf{$p$-value} \\
\midrule
Somatosensory cortex (initiation zone) & Positive & $+0.22 \pm 0.08$ & $12.4$ & $< 0.001$ \\
Thalamus (bilateral) & Positive & $+0.18 \pm 0.06$ & $10.8$ & $< 0.001$ \\
Frontal cortex & Mixed & $+0.08 \pm 0.05$ & $5.2$ & $< 0.05$ \\
Occipital cortex & Negative & $-0.12 \pm 0.05$ & $7.6$ & $< 0.01$ \\
Hippocampus & Negative & $-0.09 \pm 0.04$ & $4.8$ & $< 0.05$ \\
\bottomrule
\end{tabular}
\end{table}

\subsection{Mean-Field Simulation Reproduces Restricted Activity Propagation}

An AdEx mean-field model \citep{brette2005adaptive} was parameterized to reproduce both asynchronous irregular (AI; interictal) and SWD (ictal) dynamics by modulating the adaptation current strength. Virtual stimulation experiments confirmed that activity propagation was restricted during the ictal regime (Table~\ref{tab:simulation}).

\begin{table}[H]
\centering
\caption{Mean-field simulation results for virtual stimulation during interictal (AI) and ictal (SWD) states. Propagation extent quantified as number of connected regions showing $>$10\% activity increase. Response amplitude is peak firing rate change at the stimulated node. Values are mean $\pm$ SD across 100 simulation runs with noise.}
\label{tab:simulation}
\begin{tabular}{lcccccc}
\toprule
\textbf{State} & \textbf{Dynamics} & \textbf{Adaptation ($b$)} & \textbf{Propagation (regions)} & \textbf{Response (Hz)} & \textbf{LFP pattern} & \textbf{$p$ (AI vs SWD)} \\
\midrule
Interictal & AI & $0.02$ nS & $8.4 \pm 1.2$ & $24.6 \pm 3.8$ & Irregular & \multirow{2}{*}{$< 0.001$} \\
Ictal & SWD & $0.08$ nS & $2.1 \pm 0.8$ & $8.2 \pm 2.4$ & Spike-wave & \\
\bottomrule
\end{tabular}
\end{table}

\subsection{Stimulation Events During Scanning Sessions}

The occurrence of stimulation events classified by brain state is summarized in Table~\ref{tab:events}.

\begin{table}[H]
\centering
\caption{Stimulation events per 45-min scanning session, classified by EEG state. Values are mean $\pm$ SD across sessions. Seizure duration and count also reported.}
\label{tab:events}
\begin{tabular}{lccc}
\toprule
\textbf{Event Type} & \textbf{Visual Group ($n=6$)} & \textbf{Whisker Group ($n=5$)} & \textbf{$p$ (comparison)} \\
\midrule
Interictal stimulations & $42.3 \pm 8.4$ & $38.7 \pm 7.2$ & $0.48$ \\
Ictal stimulations & $8.6 \pm 3.2$ & $7.4 \pm 2.8$ & $0.52$ \\
Seizure-terminated-by-stim & $3.2 \pm 1.4$ & $2.8 \pm 1.2$ & $0.64$ \\
Seizures per session & $18.4 \pm 6.2$ & $16.8 \pm 5.8$ & $0.68$ \\
Mean seizure duration (s) & $4.8 \pm 1.6$ & $5.2 \pm 1.8$ & $0.72$ \\
\bottomrule
\end{tabular}
\end{table}

\section{Discussion}

This study provides the first direct evidence that absence seizures globally suppress hemodynamic responses to sensory stimulation across the entire brain. Three key findings emerge.

First, the ZTE fMRI protocol enables reliable awake-state imaging in epileptic rats with minimal motion ($<1$~$\mu$m mean displacement). The 35.8~dB noise reduction compared to EPI is critical for studying spontaneous seizures, which require a calm, unstressed awake state \citep{depaulis2015gaers}. While ZTE produces smaller BOLD signal changes ($\sim$0.3\% vs. 1--1.5\% for EPI), the reduced susceptibility artifacts and EEG gradient artifacts provide complementary advantages.

Second, both visual and somatosensory stimulation evoked positive BOLD responses during interictal periods that became negative during ictal periods. This polarity reversal suggests active suppression rather than mere absence of processing \citep{blumenfeld2005consciousness}. The suppression extended beyond primary sensory cortex to include thalamic relay nuclei and frontal attentional regions, indicating a whole-brain impairment consistent with the consciousness deficits observed clinically \citep{crunelli2020childhood}.

Third, the mean-field AdEx model provides a mechanistic account of the observed suppression \citep{brette2005adaptive, destexhe2003neuronal}. By modulating adaptation current strength, the model transitions between AI (interictal) and SWD (ictal) regimes. During SWD, the alternating burst-silent dynamics restrict the spatial propagation of stimulation-evoked activity, consistent with the reduced and spatially restricted fMRI responses observed during seizures.

The seizure hemodynamic signature---positive BOLD in somatosensory cortex (the proposed initiation zone in GAERS) and thalamus, with variable or negative responses elsewhere---is consistent with prior work on the corticothalamic origin of SWDs \citep{meeren2002cortical, polack2007deep}.

Limitations include the relatively small ZTE signal change, the inability to distinguish between hemodynamic and neural suppression with fMRI alone, and the use of a rodent model whose generalizability to human absence epilepsy, while supported by pharmacological and genetic parallels, requires verification.

\section{Methods}

\subsection{Animals and Surgery}
Eleven adult GAERS rats (8--12 months, $260 \pm 21$~g, 6 male/5 female) were implanted with carbon fiber EEG electrodes over motor and somatosensory cortex with cerebellar reference. After multi-week habituation to the restraint holder and scanner noise, animals underwent awake EEG-fMRI.

\subsection{MRI Acquisition}
A 9.4T Bruker scanner with a custom transmit-receive loop coil was used. ZTE sequence parameters were optimized for minimal acoustic noise (78.7~dB peak). Sessions lasted 45~min each.

\subsection{Stimulation Paradigm}
Visual stimulation: bilateral optical fiber cables near eyes, 6~s blocks. Whisker stimulation: bilateral air-puff tips, 6~s blocks. Stimulation was delivered during both interictal and spontaneous seizure periods; classification was performed post-hoc using EEG.

\subsection{fMRI Analysis}
GLM analysis with gamma basis functions \citep{friston1994statistical}. Voxel-wise parameter estimates computed; F-contrast maps at $p < 0.05$, cluster-corrected. ROI-based HRF analysis with extreme beta-value comparisons between states.

\subsection{Mean-Field Model}
AdEx integrate-and-fire mean-field model with excitatory and inhibitory populations. Adaptation current $b$ modulated to produce AI ($b = 0.02$~nS) or SWD ($b = 0.08$~nS) dynamics. Virtual stimulation applied as transient input current.

\subsection{Statistical Analysis}
Paired $t$-tests compared interictal vs. ictal responses within animals. Bonferroni correction for multiple ROIs. Simulation comparisons by unpaired $t$-test across runs. All tests two-sided, $\alpha = 0.05$.

\section{Conclusion}

Combining awake-state ZTE fMRI with simultaneous EEG and controlled sensory stimulation in GAERS rats, we demonstrate that absence seizures actively suppress sensory processing across the entire brain, not merely through absence of activation but through reversal of hemodynamic response polarity. Mean-field simulations provide a mechanistic account of this suppression through restricted activity propagation during ictal dynamics.

\bibliographystyle{plainnat}
\bibliography{references}

\end{document}
